\scanpage{86}\mspage{42}
Par la suite on identifie tjrs $F_M$ à une
partie de $\unsure{\mathfrak P|\mathbb P|\mathcal P}_3(S_M)$.

On \deleted{\unsure{pose}} trouve, par passage aux quotients
\deleted{(30)} à partir \deleted{\unsure{A}} de $\overrightarrow{\widetilde\tau}_F$ \deleted{$\widetilde F_M$} et $\widetilde\tau_F$ (5) (6)
des applications canoniques
\[
\begin{aligned}
\underset{\substack{\text{faces}\\\text{orientées}}}{\overrightarrow F_M}
&\xrightarrow{\overrightarrow\tau_F}\overrightarrow\pi_M=(\det M)^*\\
&&\overrightarrow\tau_F(\check f)=-\overrightarrow\tau_F(f)
\end{aligned}\tag*{(30)}
\]
\[
F_M\xrightarrow{\tau_F}\pi_M=(\det M)^*/\{\pm1_\Lambda\},\tag*{(31)}
\]
ce qui, pour $i\in\pi_M$, permet de définir
les $F_{M,i}$, d'où
\[
F_M=\coprod_{i\in\pi_M}F_{M,i}.\tag*{(32)}
\]

On pose
\[
A_M=\{a\in\mathfrak P_2(S_M)\mid\exists f\in S_M,\ f\supset a\}\tag*{(33)}
\]
et on trouve que l'on a aussi
\[
A_M=\left\{a\in\mathfrak P_2(S_M)\mid
\begin{aligned}&\exists\widetilde a\in\widetilde A_M,\ \text{avec}\\
&a=p(\widetilde a)\end{aligned}\right\}\tag*{(34)}
\]
-- mais ici le nb des $\widetilde a$ \inserted{$\in\widetilde A_M$} au dessus de $a$
est égal à 4 -- si $a=\{s,t\}$, et
$s=\dot x,t=\dot y$, les quatre choix sont
$\{x,y\},\{-x,-y\}\deleted{\unsure{\{x,y\}_M}},\{x,-y\},\{-x,y\}$
% archive sheet number: 85
